\chapter{Introduction}\label{ch:intro}

\section{The problem}

Many questions about linear dynamical systems reduce to finding the
values of a complex variable $s$ where a scalar function vanishes. The
modes of a system are the roots of its characteristic function $\Delta(s)$;
the poles and zeros of a transfer function $G(s)=N(s)/D(s)$ are the roots
of $D$ and $N$ that survive cancellation. When $\Delta$, $N$ and $D$ are
polynomials, the problem is solved by computing matrix eigenvalues, as
MATLAB's \code{roots}, \code{pole} and \code{zero} do.

Two large classes of models are not rational:
\begin{itemize}
  \item systems with \emph{time delays}, whose characteristic functions
        contain terms such as $e^{-sT}$ \citep{bellman1963,hale1993,gu2003,michiels2014};
  \item \emph{distributed-parameter} systems, described by partial
        differential equations, whose transfer functions contain
        $\sinh$, $\cosh$, $\sqrt s$ and similar functions
        \citep{curtain2009}.
\end{itemize}
Their characteristic functions typically have infinitely many roots, and
no finite matrix contains all of them. A common practice is to replace the
nonrational terms by rational approximations, for example Padé
approximations of $e^{-sT}$ or truncated modal expansions. The
approximation makes the finite-dimensional toolbox available again, but
it changes the question: one then computes the roots of a different
function. Chapter~\ref{ch:delay} shows, on simple examples, that the
answer can change qualitatively.

ContourRoots answers the original question in a bounded region: \emph{which
roots does the function have inside this rectangle?} It evaluates the
function as it is given, and uses the argument principle to count how many
roots each part of the rectangle contains, which makes it possible to tell
whether the search was complete.

\section{What the toolbox provides}

\begin{itemize}
  \item \code{croots}, \code{cpoles}, \code{czeros} and \code{cpzmap}, a
        MATLAB-style interface for roots, poles, zeros and pole-zero maps
        of scalar functions given as function handles, coefficient vectors,
        numerator/denominator pairs (\code{ndpair}), symbolic expressions or
        SISO LTI models.
  \item Detection of pole-zero cancellations, including partial ones.
  \item Diagnostics for every search: a status, the multiplicity and
        residual of each location, cancelled locations and unresolved
        regions.
  \item For scalar retarded time-delay systems $D(s)+N(s)e^{-sT}=0$: all
        critical delays and crossing directions up to a given delay,
        without approximation, and a count of unstable roots in a region
        that provably contains all of them.
  \item Examples, reproducible studies and tests.
\end{itemize}

\section{What the toolbox does not claim}

The word \emph{exact} is used in this manual only in one sense: the original
function is evaluated, without rational approximation or truncation. The
computed roots are floating-point numbers. A search that ends with the
status \code{numerically\_complete} has passed careful numerical checks
under the analyticity assumption, but it is not a proof in the sense of
interval arithmetic (Section~\ref{sec:contract}). The toolbox does not treat
matrix-valued characteristic functions, does not enumerate infinite spectra,
and does not discover singularities hidden inside opaque function handles.
No claim of superiority over other root finders is made; related methods
and software are discussed in Section~\ref{sec:related}.

\section{How to read this manual}

Chapter~\ref{ch:quick} is a short tour of the interface.
Chapters~\ref{ch:background}--\ref{ch:inputs} give the mathematical
background, the algorithm, and the precise meaning of each input form and
of the analyticity assumption. Chapter~\ref{ch:api} is a compact reference.
Chapters~\ref{ch:delay}--\ref{ch:beam} are case studies that can be read
independently: time-delay systems and the limits of Padé approximations;
distributed-parameter systems; and a beam coupled to an oscillator.
Chapter~\ref{ch:validation} summarizes the validation, and the appendices
contain proofs and reproducibility information.

The Markdown documentation distributed with the toolbox covers the same
material in a tutorial style, with runnable code for every step.
